By the score-expansion and covariance-consistency lemmas,
\[
 \sqrt n\{\widehat\mu-B\Theta(P)\}
 =\frac1{\sqrt n}\sum_{i=1}^n\Psi(O_i)+o_P(1),
 \qquad
 \widehat\Omega_\lambda\to_P\Omega
\]
in fixed dimension.  The spectrum assumption makes \(\Omega\) positive definite.  The incidence matrix \(B\) has full column rank because every row belongs to exactly one nonempty block.  Matrix inversion is continuous on the positive-definite cone, so
\[
 \widehat b^{\mathrm{joint}}
 =\widehat\Omega_\lambda^{-1}B
   (B'\widehat\Omega_\lambda^{-1}B)^{-1}\omega
 \to_P
 \Omega^{-1}B(B'\Omega^{-1}B)^{-1}\omega
 =b^{\mathrm{joint}}.
\]
Both the estimated and population loadings satisfy \(B'b=\omega\).  Consequently,
\[
\begin{aligned}
 \sqrt n\{\widehat\tau_{\mathrm{joint}}-\Theta_\omega^{\mathrm{agg}}(P)\}
 &=\widehat b^{\mathrm{joint}\prime}
   \sqrt n\{\widehat\mu-B\Theta(P)\}\\
 &=\frac1{\sqrt n}\sum_{i=1}^n
   b^{\mathrm{joint}\prime}\Psi(O_i)+o_P(1).
\end{aligned}
\]
Indeed, the empirical influence sum is \(O_P(1)\) by its bounded second moment, so replacing \(\widehat b^{\mathrm{joint}}\) by its probability limit costs \(o_P(1)\).

Set \(Z_i=b^{\mathrm{joint}\prime}\Psi(O_i)\).  These variables are i.i.d., have mean zero, variance \(V_{\mathrm{joint}}\), and a finite fourth moment.  To record the scalar central-limit step explicitly, its characteristic function satisfies
\[
 E_P[e^{iuZ_i/\sqrt n}]
 =1-\frac{u^2V_{\mathrm{joint}}}{2n}+o(n^{-1}),
\]
where the remainder follows by truncation and the finite second moment.  Independence and passage to the \(n\)-th power give the limiting characteristic function
\(\exp\{-u^2V_{\mathrm{joint}}/2\}\), hence
\[
 \frac1{\sqrt n}\sum_{i=1}^nZ_i\rightsquigarrow N(0,V_{\mathrm{joint}}).
\]

We next establish optimality.  For the strictly convex program
\[
 \min_{b:\,B'b=\omega} b'\Omega b,
\]
the first-order equations are \(\Omega b=B\lambda\) and \(B'b=\omega\).  Solving gives
\(b=\Omega^{-1}B(B'\Omega^{-1}B)^{-1}\omega=b^{\mathrm{joint}}\).  Strict convexity makes this minimizer unique.  It is nonzero because \(B'b^{\mathrm{joint}}=\omega\) and \(\mathbf 1'\omega=1\).  Therefore the spectrum bound yields
\[
 V_{\mathrm{joint}}
 =b^{\mathrm{joint}\prime}\Omega b^{\mathrm{joint}}
 \ge\kappa\|b^{\mathrm{joint}}\|^2>0.
\]
Covariance consistency and loading consistency give
\[
 \widehat V_{\mathrm{joint}}
 =\widehat b^{\mathrm{joint}\prime}\widehat\Omega
   \widehat b^{\mathrm{joint}}
 \to_PV_{\mathrm{joint}}.
\]
The identification lemma equates \(\Theta_\omega^{\mathrm{agg}}(P)\) with \(\tau_\omega(Q)\).  Dividing the linear expansion by \(\sqrt{\widehat V_{\mathrm{joint}}}\) and using the two probability limits above yields a standard normal limit.  Since the standard normal distribution is continuous at \(\pm z_{1-\alpha/2}\), the event defining the Wald interval has limiting probability \(1-\alpha\).  Positivity of the limiting variance also shows that the finite-sample zero-variance convention is asymptotically inactive.

If \(g=2\), the definition of \(\Delta Y_{i,k}\) gives \(Y_{it}-Y_{i1}\) in both oracle and fitted scores.  Finally, if \(\mathcal K=\{k\}\), then \(B=\mathbf 1\) and the constraint forces \(\omega_k=1\); hence
\[
 b^{\mathrm{joint}}=\frac{\Omega^{-1}\mathbf 1}{\mathbf 1'\Omega^{-1}\mathbf 1},
\]
which is exactly the componentwise comparison-cohort loading, with the same feasible estimator and variance limit.